\chapter{Hermon DNA prefill, decode and KV state}
\label{evod35:chapter}

A Transformer can recompute the entire consumed prefix before predicting
each next token. That is a useful correctness oracle and a costly serving
strategy. A key/value cache retains intermediate results that future queries
need, avoiding repeated old-token projections. But the cached numbers must
come from the right model, layer, position convention and causal prefix.
A tensor with the right shape can still be the wrong state.

This chapter implements caching for the \emph{actual Chapter 23 Hermon DNA
reference}: a pre-normalized Transformer with learned
absolute positions, layer normalization, attention, residuals, GELU
feed-forward sublayers, final normalization and an ACGT language-model head.
We do not silently replace it with a single-head rotary-attention toy.
The restricted executable fixture uses two layers, two heads, width eight,
CPU float64 and deterministic evaluation. All parameters are assigned by a
seeded initialization; no trained genomic capability is asserted.

We will first identify what persists, derive prefill and offset-masked chunk
continuation, prove equivalence through the complete stack, and then count
payload and copying costs. Chapter 36 may change allocation, but not these
logical semantics.

\section{Specify which Hermon computation is being served}
\label{evod35:scope}
Chapter 23 supports causal logits and offline full-record classification.
Only its \emph{causal} language-model path admits the ordinary left-to-right
cache developed here. Offline attention or reverse-complement averaging can
depend on unseen record suffixes. Their state is not interchangeable with a
causal cache.

A request is a single unpadded canonical ACGT token record
$u_0,\ldots,u_{t-1}$; the mapping is $A=0,C=1,G=2,T=3$.
Ambiguous bases, k-mer spans, multiple packed records and padded ragged batches
are outside this adapter. Chapter 23's validity contract remains available
in the full model, but the serving adapter does not pretend to implement all
of it. A new record starts with an empty cache; carrying another record's
history changes the conditional distribution.

The original model forms
\[
 x_i^0=E_{u_i}+P_i,
\]
where $P_i$ is a learned position embedding. Position is the consumed-token
ordinal in this unpadded contract. For variable-span genomic tokens it would
not automatically equal a base coordinate; that mapping must come from the
data/coordinate contract, not a kernel's array index.

\section{Trace one block and locate its retained state}
\label{evod35:block}
Let $x^\ell_i\in\mathbb R^D$ be the input row of block $\ell$. With $H$
heads and $d_h=D/H$, the block computes
\begin{align*}
 r^\ell_i&=\operatorname{LN}_1^\ell(x^\ell_i),\\
 (q^\ell_i,k^\ell_i,v^\ell_i)&=
       \operatorname{QKV}^\ell(r^\ell_i),\\
 y^\ell_i&=x^\ell_i+
       W_O^\ell\operatorname{MergeHeads}(\operatorname{Attention}_i^\ell)
       +b_O^\ell,\\
 x^{\ell+1}_i&=y^\ell_i+
       \operatorname{FFN}^\ell(\operatorname{LN}_2^\ell(y^\ell_i)).
\end{align*}
The displayed output projection uses the mathematical column-vector
convention. The reference uses PyTorch's linear layers and row tensors;
it does not transpose or discard the learned biases.

Each key and value is computed from the \emph{current layer's normalized
activation}, not directly from the original token embedding. At layer two,
that activation already includes layer one's causal attention and residual
processing. Caching only first-layer token projections and reusing them at
every layer is therefore wrong, even when dimensions happen to match.

\Plate{02-geometry}{Query rows are temporary. Ordered K/V rows persist separately at each normalized attention projection.}{evod35:geometryplate}

For one head,
\[
 s_{ij}=\frac{\langle q_i,k_j\rangle}{\sqrt{d_h}},
 \qquad
 a_{ij}=\frac{\exp s_{ij}}{\sum_{r\le i}\exp s_{ir}},
 \qquad o_i=\sum_{j\le i}a_{ij}v_j.
\]
Softmax is taken over allowed keys, not over heads, channels or requests.
Use max subtraction for stable evaluation; the library softmax in the
reference supplies that operation. Every query in this contract has at least
its own current key, so it never needs an all-masked-row fallback.

The scaled dot-product and multi-head construction follows the Transformer
formulation~\cite{vaswani}. The particular pre-normalized architecture and
learned-position choice belong to our existing reference, not to a claim that
every Transformer has these exact conventions.

\subsection{A small retrieval calculation}
Suppose one query has three allowed scores $(0,\log2,0)$ and values
$(1,0),(0,2),(2,1)$. The normalized weights are $(1/4,1/2,1/4)$, giving
\[
 o=(3/4,5/4).
\]
This calculation clarifies why caching old keys and values is enough:
the new query needs the old \emph{addresses and contents}, not the old queries.
The chapter tests this assigned retrieval independently of the model fixture.

\section{Define the cache before optimizing it}
\label{evod35:cache}
A \Term{per-layer KV cache} at consumed length $t$ contains
\[
 C_t^\ell=(K_t^\ell,V_t^\ell),\qquad
 K_t^\ell,V_t^\ell\in\mathbb R^{H\times t\times d_h}.
\]
This is retained addressable history. It is neither Chapter 31's compressed
recurrent vector nor the Transformer parameters. A cache's position is the
number of consumed tokens, not the number of generated tokens since prefill.

The envelope binds owner, artifact, exact causal mode, ordered consumed-token
metadata and the ordered collection of layer caches. Artifact identity includes
the actual model's parameter tensors, position table, geometry,
normalization epsilon settings and GELU approximation convention.
The reference requires evaluation-mode finite CPU float64 parameters and
bounded geometry before it accepts a cache.

\Plate{04-context}{Cache correctness has context, geometry and integrity layers. None alone proves that stored values were produced by native replay.}{evod35:contextplate}

The dataclasses are frozen, but tensors inside them are not intrinsically
immutable. The adapter therefore creates fresh storage on continuation and
checks a \Term{cache content seal} over metadata and tensor values before
use. This detects accidental mutation. A producer with access to both content
and hashing can alter values and reseal them; the negative control demonstrates
that such a structurally valid cache can change logits.

The seal is not authentication, a cryptographic authorization capability,
a proof of correct prefix computation, hostile-process isolation or an
untrusted-file loader. Nor does storing token IDs prove the tensors were
derived from those IDs. Trusted producers and serialized callers are explicit
assumptions.

\Needspace{29\baselineskip}
\Listing{validate}

Validation deliberately rejects dtype changes rather than silently converting
the history. Casting a cache after construction changes the numerical
contract; quantized cache designs require their own parity and quality
experiments. Content hashing and tuple conversion are expensive teaching
checks, not a proposed per-token production fast path.

\section{Prefill constructs a causal history}
\label{evod35:prefill}
\Term{Prefill} consumes a prompt of $T$ tokens starting from an empty cache.
All prompt positions can be projected together, but query $i$ can read only
keys $j\le i$. Parallel evaluation is not bidirectional information flow.
Each layer stores its full prompt K/V, and the final normalized rows pass
through the language-model head.

The last prompt row's logits predict the token \emph{after} the prompt. When
that selected token is later consumed by decode, its new logits predict the
following token. Confusing predicted with consumed positions produces an
off-by-one error that a shape check cannot detect. This chapter uses supplied
tokens rather than a sampling loop to isolate execution parity.

\Plate{01-equivalence}{Different evaluation schedules must preserve all rows of the original full causal computation, not merely the final cache length.}{evod35:equivalenceplate}

The simplest independent oracle is the model's original
\texttt{causal\_logits} method. It performs complete full-prefix computation
with its own square blocked-mask construction and never calls the cache
append helper.
\Listing{full_oracle}

Calling prefill twice through the same cache helper would not constitute
an independent full-versus-cached oracle: the same offset or layer-order bug
could appear on both sides.

\section{Decode and causal chunks need the correct offset}
\label{evod35:offset}
Let $t$ old tokens already be cached and let a continuation chunk contain $c$
new tokens. New activations have shape $[c,D]$ and positions
$t,\ldots,t+c-1$. At each layer,
\[
 Q_{\rm new}\in\mathbb R^{H\times c\times d_h},\qquad
 K_{\rm all},V_{\rm all}\in\mathbb R^{H\times(t+c)\times d_h}.
\]
Old rows and current rows concatenate in causal order. Local query index $i$
is absolute position $t+i$, so its allowed-key condition is
\[
 A_{ij}=[j\le t+i],\quad
 0\le i<c,\quad0\le j<t+c.
\]
The old prefix is entirely visible; only the future portion of the new chunk
must be masked.

\Plate{03-offset}{A rectangular continuation mask uses absolute query positions. An upper-left triangle can be shape-correct and semantically wrong.}{evod35:offsetplate}
\Listing{allowed_mask}

The reference's mask uses \texttt{True} to mean \emph{allowed}. Chapter 23's
\texttt{legality} helper uses \texttt{True} to mean \emph{blocked}.
The conversion in \texttt{consume} is explicit. PyTorch 2.10 SDPA accepts an
allowed boolean mask; its default non-square causal alignment is upper-left,
which is not this offset continuation mask~\cite{sdpa}. Tests compare an
explicit-mask SDPA oracle and deliberately show that the upper-left shortcut
changes the result for a two-query, five-key continuation.

For \Term{one-token decode}, $c=1$ and all $t+1$ stored/current keys are allowed.
No triangular mask is needed then, \emph{provided} the cache really contains
only this causal prefix. The more general explicit-mask implementation covers
both one-token decode and multi-token chunk continuation.

\subsection{The complete append path}
\Listing{consume}

A continuation call validates old state and the entire token chunk before
arithmetic. It constructs candidates in new storage, checks finite logits
and cache payloads, then returns a new cache. The caller decides whether to
install it as the current request state. The old cache remains unchanged on
success and on validation/numerical failure. This is functional candidate
construction, not a thread-safe mutable cache transaction.

The special empty chunk returns zero rows and the exact same cache object.
It is an identity, not an implicit BOS token. An empty prefill likewise
produces no logits: a caller cannot select a next token from a nonexistent
last row without defining a separate beginning-of-sequence convention.

\section{Prove parity through tokens and layers}
\label{evod35:proof}
\Needspace{14\baselineskip}
\begin{proposition}[Full/cached causal equivalence]
Fix model parameters and architecture, deterministic evaluation, token IDs,
causal mode, position convention and record identity. Suppose a cache contains
the exact per-layer K/V produced by full causal evaluation of the consumed
prefix. Appending a valid continuation chunk computes the same new output
rows and extended layer caches as full causal evaluation of the combined
record, in real arithmetic.
\end{proposition}
\begin{proof}
At layer zero, each new row receives the same token and absolute position
embedding as full evaluation. Assume the old cached rows match full K/V and
the new input activations match at the current layer. Normalization is
row-local, so the new Q/K/V projections match. The offset mask selects exactly
the same keys as the corresponding full causal row, including old prefix and
allowed current-chunk rows. Softmax and value aggregation therefore match.
The output projection, residual additions and feed-forward sublayer are
row-local and use identical parameters, so the next layer's new inputs match.
Induction over layers proves matching logits and extended K/V. Induction over
chunks preserves the hypothesis for any continuation partition.
\end{proof}

Causality is essential to the old-cache hypothesis: appending future tokens
cannot change old causal activations. A bidirectional layer would invalidate
it. Deterministic evaluation is also essential to this argument as stated.
The reference has no dropout modules; a future backend that introduces
attention dropout must pass an explicitly zero inference dropout probability,
not assume a surrounding evaluation flag overrides the functional API.

In floating-point arithmetic, full and chunk paths may use differently
shaped matrix reductions. Tests use absolute and relative tolerances
$2\times10^{-12}$ in CPU float64. This is a checked fixture tolerance,
not a universal backward-error theorem for every model, precision or device.
Per-layer key/value checks narrow the location of a failure; final-logit
parity alone can hide compensating mistakes.

\begin{center}\input{\Generated/parity.tex}\end{center}
The table is generated from the two-layer six-token fixture. Every prefill
split is followed by scalar decode through the actual stack. All end at six
consumed tokens and 1,536 logical KV bytes. Separate tests cover every one
of the $2^5=32$ nonempty chunk compositions of that record.

\section{Positions cannot be repaired by tensor shape}
\label{evod35:positions}
Our reference's learned absolute embeddings enter \emph{before} layer one.
Restarting their ordinal at zero after prefill changes Q/K/V throughout the
stack. Unlike standard rotary relative-score identities, learned position
tables do not in general permit a uniform coordinate shift with invariant
outputs. Importing such an identity from another architecture would be wrong.

A cache artifact therefore binds the position table and model semantics,
and the next position is derived from consumed length. A negative control
copies position-zero embedding into the next-position entry of a separately
identified model: earlier prefix state stays comparable, but continuation
logits change. Normal context validation would reject a mid-request artifact
switch before computing those wrong logits.

Changing strand orientation also changes the ordered token sequence and
positions. A forward causal cache cannot simply be relabelled as a
reverse-complement cache. An offline symmetry construction from earlier
chapters is not a legal hidden transformation of a left-to-right serving
state. A transformed-cache reuse rule would need a separate proof for the
chosen model and readout.

\section{Count resident payload separately from copying}
\label{evod35:memory}
For $L$ layers, $t$ retained tokens, $H_{kv}$ key/value heads, head width
$d_h$ and element width $b$ bytes,
\[
 M_{\rm KV}(t)=2LtH_{kv}d_hb.
\]
The factor two counts K and V. Query heads do not persist. With ordinary
multi-head attention, $H_{kv}=H_q=H$; grouped-query and multi-query designs
use fewer KV heads~\cite{gqa}. That changes the architecture and cache
geometry. Our executable adapter implements the existing equal-head model,
not a conversion to GQA, and asserts no quality-preserving head reduction.

\Plate{05-cost}{Resident history grows linearly. Copying old contiguous prefixes on every append creates an additional, cumulative growth cost.}{evod35:costplate}
\begin{center}\input{\Generated/memory.tex}\end{center}

The large arithmetic example uses $L=24,H_{kv}=8,d_h=128,b=2$:
at $t=4096$, $M=402{,}653{,}184$ bytes or 384 MiB for one request.
These dimensions describe a storage calculation, not the small executable
fixture or a trained Hermon artifact. Parameters, activations, metadata,
stored token IDs, allocator reserve and temporary copies are excluded.
Actual device allocation must be measured separately.

\subsection{The cost hidden by a simple concatenate}
Define $\gamma=2LH_{kv}d_hb$, the logical KV bytes per token. Starting with
prompt length $P$, append $G$ tokens one at a time. An implementation that
copies the old prefix into a fresh destination performs old-prefix reads and
equal-length writes totaling
\[
 C_{\rm old}
 =2\gamma\sum_{g=0}^{G-1}(P+g)
 =2\gamma GP+\gamma G(G-1).
\]
This term is quadratic in generated length even though final resident payload
$\gamma(P+G)$ is linear. New KV-row writes add $\gamma G$; they are not
included in $C_{\rm old}$.

For the executable geometry, $\gamma=256$ bytes. With $P=3,G=3$, the
old-prefix read/write component is $2(256)(3+4+5)=6144$ bytes. This is not
the reference's measured total memory traffic: it additionally clones new
layer tensors, creates temporary projections and hashes values. The formula
isolates why repeated relocation becomes undesirable rather than pretending
the Python code has achieved the lower bound.

\section{Work, traffic and the limits of a serving comparison}
\label{evod35:work}
A dense prefill of $T$ tokens forms a $T\times T$ score array per head.
Its mathematical causal dependency graph has $T(T+1)/2$ allowed pairs, but
the transparent implementation still computes dense score and value
products. A triangular dependency diagram is not proof that the code
skips half the arithmetic.

For a new query at retained length $t$, score and value interactions require
work proportional to $td_h$ per query head; projection and feed-forward
work remain. The amount of logically retained history that may need reading
also grows with $t$. Whether decode is bandwidth-limited on particular hardware
depends on batch size, cache reuse, dtype, weight traffic and kernel layout;
it is a profile question, not an unconditional law.

FlashAttention demonstrates an IO-aware organization of exact attention
that avoids materializing the full attention matrix in high-bandwidth
memory~\cite{flash}. This is a relevant algorithmic comparison. It does not
mean this chapter's dense CPU function is FlashAttention or that finite-precision
outputs from every backend are bitwise equal. Chapter 38 owns optimized
attention-kernel experiments.

Comparing the fixed-width recurrence in Chapter 34 with this growing KV
representation establishes a \emph{state-geometry distinction}. It does not
establish equal expressive capacity, genomic accuracy, latency, energy
efficiency or a universal winner. Those require shared tasks, budgets,
hardware and appropriate evidence from Chapters 25 and 29.

\section{Verification and the allocator boundary}
\label{evod35:verification}
The local reference test suite checks every fixture prefill split across
three model seeds, every chunk composition, original full-model QKV hooks at
each layer, future-prefix invariance, fresh storage, empty identity,
capacity rejection, model/position/configuration changes, invalid tokens,
owner/mode mismatches, finite payload mutation, nonfinite payloads, wrong
geometry/dtype/gradient state, and a resealed forgery.

A content seal can report a mutated tensor before use; it cannot explain
whether a trusted producer supplied a false cache originally. Full native
replay could check consistency but would itself cost prefix computation.
The chapter neither authenticates arbitrary incoming cache files nor calls
an executable object deserializer.

\Plate{06-refinement}{Chapter 36 may relocate K/V physically while preserving this chapter's exact logical history and ownership.}{evod35:refinementplate}

Preallocating a bounded contiguous buffer can avoid repeated relocation,
but it creates capacity reservation and mutable-lifetime questions. Paging
can address a different trade-off, allowing noncontiguous physical storage
and prefix sharing. These are allocator refinements. They must preserve the
ordered per-layer history visible to each query, artifact and coordinate
identity, and reference lifetime. They do not alter the attention formula
to excuse a missing or duplicated key.

\section{Twenty worked exercises}
\label{evod35:exercises}
\begin{enumerate}
\item \textbf{Find the cache boundary.} Why cache normalized projected K/V,
not token embeddings alone? \emph{Solution:} later-layer K/V depend on
earlier-layer attention, residuals and FFN. Embeddings alone would require
recomputing that causal stack to recover them.
\item \textbf{Normalize three scores.} Evaluate $(0,\log2,0)$.
\emph{Solution:} exponentials are $(1,2,1)$; divide by four. The assigned
values then aggregate to $(3/4,5/4)$.
\item \textbf{Explain transient queries.} Why not retain old Q?
\emph{Solution:} future causal rows use new queries against old K/V. Old
outputs are not recomputed in ordinary append-only decode.
\item \textbf{Derive offset legality.} For past length three and chunk two,
which keys can the first new query read? \emph{Solution:} absolute position
three reads indices zero through three, not four.
\item \textbf{Expose the rectangular shortcut.} What does an upper-left
$2\times5$ triangle do? \emph{Solution:} its rows read only one and two keys.
The correct offset mask reads four and five.
\item \textbf{Convert boolean conventions.} Chapter 23 returns blocked
edges; SDPA accepts allowed edges. \emph{Solution:} invert the boolean mask
explicitly. Matching shapes do not reconcile opposite meanings.
\item \textbf{Locate the next token.} What do the last prompt logits predict?
\emph{Solution:} the token after the consumed prompt. Consuming that token
then advances cache length and predicts another token.
\item \textbf{Handle empty input.} Does empty prefill supply a next-token
distribution? \emph{Solution:} no, it produces no output row. Define an
explicit start-token policy if a generation API needs one.
\item \textbf{Prove layer induction.} What is the essential inductive
hypothesis? \emph{Solution:} old per-layer K/V and current new-layer inputs
equal their full causal counterparts; row-local operations preserve equality.
\item \textbf{Count partitions.} Why 32 compositions for six tokens?
\emph{Solution:} each of the five internal gaps independently contains a
chunk boundary or not.
\item \textbf{Test the whole model.} Is a single-head attention comparison
enough? \emph{Solution:} no; include layer-wise projections, residuals,
normalizations, FFN, position embeddings and the final head.
\item \textbf{Shift absolute positions.} Does a uniform shift preserve this
model? \emph{Solution:} not generally; learned position rows are unrelated
parameters. Do not import a rotary relative-position identity.
\item \textbf{Switch strands.} Can forward KV become reverse-complement KV
by changing metadata? \emph{Solution:} no. The sequence and causal positions
change; a valid transformation would need a model-specific proof.
\item \textbf{Compute fixture payload.} Two layers, six tokens, two heads,
head width four and eight-byte elements? \emph{Solution:}
$2(2)(6)(2)(4)(8)=1536$ bytes, excluding all overhead.
\item \textbf{Compute the large example.} Convert 402,653,184 bytes to MiB.
\emph{Solution:} divide by $2^{20}$ to obtain 384 MiB, not 402.65 MiB.
\item \textbf{Separate head geometries.} Which head count controls KV?
\emph{Solution:} $H_{kv}$; query heads control a different part of attention
work. Changing KV heads is not merely an allocator optimization.
\item \textbf{Derive copy growth.} Why quadratic cumulative traffic?
\emph{Solution:} each append recopies a prefix one token longer; summing
$P,\ldots,P+G-1$ creates $GP+G(G-1)/2$.
\item \textbf{Detect versus authenticate.} Can a producer reseal altered
K/V? \emph{Solution:} yes. Structural integrity accepts it, but a native
replay/logit negative control exposes inconsistency.
\item \textbf{Describe a numerical guarantee honestly.} Does fixture
float64 parity certify quantized GPU KV? \emph{Solution:} no; dtype,
kernel reductions, error propagation and downstream quality need new tests.
\item \textbf{State the paging invariant.} What must a page table preserve?
\emph{Solution:} exactly ordered layer/head/token K/V values, model and
position identity, owner isolation and lifetime. Physical adjacency is not
part of the attention semantics.
\end{enumerate}

\section{Handoff: retain the oracle while changing storage}
Chapter 36 should compare paged allocation against this exact causal oracle,
including partial final blocks, branch divergence, reference counting and
reclamation. Eviction may remove a cache's ownership reference; it must not
invalidate pages still used by a live request. Establishing those invariants
will turn a storage diagram into an executable system argument without
changing what the model computes.
